\documentclass[fleqn]{ctexart}
\usepackage{graphicx}
\usepackage{xcolor}
% 数学支持（提供 \mathbb、\dfrac 等）
\usepackage{amsmath,amssymb}
% 页面与版式设置：更小的页边距与顶部空白
\usepackage[a4paper,top=15mm,bottom=20mm,left=12mm,right=12mm]{geometry}
% 自动换行与字偶距优化
\usepackage{microtype}
\microtypesetup{tracking=true,protrusion=true,final}
% 在需要时放宽断行以避免溢出（数值可按需调整）
\setlength{\emergencystretch}{2em}
% 列表控制，便于让(1)后直接跟内容
\usepackage{enumitem}
% verbatim 环境支持（用于直接粘贴源码与输出）
\usepackage{verbatim}
\usepackage{fancyvrb}
\usepackage{fvextra}
\RecustomVerbatimEnvironment{verbatim}{Verbatim}{
	frame=single,
	framesep=2mm,
	breaklines=true,
	breakanywhere=true
}
% 使章节标题左对齐（同时保留 ctex 默认样式）
\ctexset{section={format+=\raggedright}}
%1.3倍行距
\renewcommand{\baselinestretch}{1.3}
% 数学左对齐缩进为0
\setlength{\mathindent}{0pt}
% 增大矩阵的行距
\renewcommand{\arraystretch}{1.3}
\pagestyle{empty}
\begin{document}

\section*{题目1：空指针解引用（NULL）}
\subsection*{源代码}
\begin{verbatim}
#include <stdio.h>

int main(void)
{
	int *p = NULL;
	printf("p=%p\n", (void *)p);
	*p = 42;
	printf("unreachable\n");
	return 0;
}
\end{verbatim}

\subsection*{直接运行现象}
\begin{verbatim}
p=(nil)
exit:139
\end{verbatim}

\subsection*{分析}
该程序将指针 \verb|p| 置为 \verb|NULL|（地址 0），随后执行 \verb|*p = 42|。
这会向地址 0 写入一个 \verb|int|（4 字节），属于空指针解引用引发的非法写内存。

CPU 触发保护异常，进程收到 \verb|SIGSEGV| 并终止。
本次运行的退出码为 139，通常对应 128+11，即被信号 11（\verb|SIGSEGV|）杀死。


\section*{题目2：使用 gdb 调试空指针解引用}
\subsection*{源代码}
\begin{verbatim}
#include <stdio.h>

int main(void)
{
	int *p = NULL;
	printf("p=%p\n", (void *)p);
	*p = 42;
	printf("unreachable\n");
	return 0;
}
\end{verbatim}

\subsection*{gdb 输出}
\begin{verbatim}
[Thread debugging using libthread_db enabled]
Using host libthread_db library "/lib/x86_64-linux-gnu/libthread_db.so.1".
p=(nil)

Program received signal SIGSEGV, Segmentation fault.
0x000055555555519c in main () at 02_null_gdb.c:7
7           *p = 42;
#0  0x000055555555519c in main () at 02_null_gdb.c:7
#0  0x000055555555519c in main () at 02_null_gdb.c:7
7           *p = 42;
p = 0x0
$1 = (int *) 0x0
\end{verbatim}

\subsection*{分析}
gdb 显示程序收到了 \verb|SIGSEGV|，并停在 \verb|02_null_gdb.c:7| 的 \verb|*p = 42| 这一行。
这说明崩溃点就是对空指针的写操作。
回溯（\verb|bt|）只有 \verb|main| 一帧，符合该示例程序结构。
同时 gdb 打印 \verb|p = 0x0|，确认指针值为 0（NULL），因此对其解引用必然触发非法访问。


\section*{题目3：使用 valgrind 分析空指针解引用}
\subsection*{源代码}
\begin{verbatim}
#include <stdio.h>

int main(void)
{
	int *p = NULL;
	printf("p=%p\n", (void *)p);
	*p = 42;
	printf("unreachable\n");
	return 0;
}
\end{verbatim}

\subsection*{valgrind 输出}
\begin{verbatim}
==138968== Memcheck, a memory error detector
==138968== Copyright (C) 2002-2022, and GNU GPL'd, by Julian Seward et al.
==138968== Using Valgrind-3.22.0 and LibVEX; rerun with -h for copyright info
==138968== Command: ./03_null_valgrind
==138968== 
p=(nil)
==138968== Invalid write of size 4
==138968==    at 0x10919C: main (03_null_valgrind.c:7)
==138968==  Address 0x0 is not stack'd, malloc'd or (recently) free'd
==138968== 
==138968== 
==138968== Process terminating with default action of signal 11 (SIGSEGV)
==138968==  Access not within mapped region at address 0x0
==138968==    at 0x10919C: main (03_null_valgrind.c:7)
==138968==  If you believe this happened as a result of a stack
==138968==  overflow in your program's main thread (unlikely but
==138968==  possible), you can try to increase the size of the
==138968==  main thread stack using the --main-stacksize= flag.
==138968==  The main thread stack size used in this run was 8388608.
==138968== 
==138968== HEAP SUMMARY:
==138968==     in use at exit: 1,024 bytes in 1 blocks
==138968==   total heap usage: 1 allocs, 0 frees, 1,024 bytes allocated
==138968== 
==138968== LEAK SUMMARY:
==138968==    definitely lost: 0 bytes in 0 blocks
==138968==    indirectly lost: 0 bytes in 0 blocks
==138968==      possibly lost: 0 bytes in 0 blocks
==138968==    still reachable: 1,024 bytes in 1 blocks
==138968==         suppressed: 0 bytes in 0 blocks
==138968== Reachable blocks (those to which a pointer was found) are not shown.
==138968== To see them, rerun with: --leak-check=full --show-leak-kinds=all
==138968== 
==138968== For lists of detected and suppressed errors, rerun with: -s
==138968== ERROR SUMMARY: 1 errors from 1 contexts (suppressed: 0 from 0)
Segmentation fault (core dumped)
\end{verbatim}

\subsection*{分析}
valgrind 的 \verb|Invalid write of size 4| 表示发生了一次非法写入，写入大小为 4 字节（与 \verb|int| 大小一致）。
报告定位到 \verb|03_null_valgrind.c:7|，对应语句 \verb|*p = 42|。

\verb|Address 0x0 is not stack'd, malloc'd or (recently) free'd| 说明目标地址 0x0 不属于栈区、堆区（malloc/realloc 分配的区域），也不是刚释放的块，因此属于无效地址访问。

后续的 \verb|SIGSEGV| 与“\verb|Access not within mapped region at address 0x0|”进一步确认是对未映射地址的访问。
此外 \verb|HEAP/LEAK SUMMARY| 里的 \verb|still reachable| 多为运行时库内部分配（例如 I/O 缓冲），关键问题仍是非法写导致崩溃。


\section*{题目4：malloc 后忘记 free（内存泄漏）}
\subsection*{源代码}
\begin{verbatim}
#include <stdio.h>
#include <stdlib.h>

int main(void)
{
	size_t n = 1024;
	char *buf = (char *)malloc(n);
	if (buf == NULL)
	{
		perror("malloc");
		return 1;
	}

	buf[0] = 'A';
	buf[n - 1] = 'Z';
	printf("buf[0]=%c, buf[%zu]=%c\n", buf[0], n - 1, buf[n - 1]);

	return 0;
}
\end{verbatim}

\subsection*{直接运行现象}
\begin{verbatim}
buf[0]=A, buf[1023]=Z
exit:0
\end{verbatim}

\subsection*{valgrind 输出}
\begin{verbatim}
==139136== Memcheck, a memory error detector
==139136== Copyright (C) 2002-2022, and GNU GPL'd, by Julian Seward et al.
==139136== Using Valgrind-3.22.0 and LibVEX; rerun with -h for copyright info
==139136== Command: ./04_leak
==139136== 
buf[0]=A, buf[1023]=Z
==139136== 
==139136== HEAP SUMMARY:
==139136==     in use at exit: 1,024 bytes in 1 blocks
==139136==   total heap usage: 2 allocs, 1 frees, 2,048 bytes allocated
==139136== 
==139136== 1,024 bytes in 1 blocks are definitely lost in loss record 1 of 1
==139136==    at 0x4846828: malloc (in /home/amadeus/.local/valgrind/usr/libexec/valgrind/vgpreload_memcheck-amd64-linux.so)
==139136==    by 0x1091A8: main (04_leak.c:7)
==139136== 
==139136== LEAK SUMMARY:
==139136==    definitely lost: 1,024 bytes in 1 blocks
==139136==    indirectly lost: 0 bytes in 0 blocks
==139136==      possibly lost: 0 bytes in 0 blocks
==139136==    still reachable: 0 bytes in 0 blocks
==139136==         suppressed: 0 bytes in 0 blocks
==139136== 
==139136== For lists of detected and suppressed errors, rerun with: -s
==139136== ERROR SUMMARY: 1 errors from 1 contexts (suppressed: 0 from 0)
\end{verbatim}

\subsection*{分析}
该程序分配了 1024 字节（\verb|malloc(n)|）并正常使用，但退出前没有 \verb|free(buf)|。
由于进程结束时操作系统会回收该进程的整个虚拟地址空间，所以直接运行看起来“完全正常”。

valgrind 会在进程退出时做堆内存审计：
\begin{itemize}[leftmargin=2em]
\item \verb|in use at exit|：退出时仍未释放的堆内存。
\item \verb|definitely lost|：确定泄漏。
\item 本例显示 \verb|1,024 bytes ... definitely lost|，并回溯到 \verb|04_leak.c:7| 的 \verb|malloc|。
\end{itemize}

结论：程序功能正确，但存在内存泄漏，需要在返回前添加 \verb|free(buf)|。


\section*{题目5：数组越界写（data[100]）}
\subsection*{源代码}
\begin{verbatim}
#include <stdio.h>
#include <stdlib.h>

int main(void)
{
	int *data = (int *)malloc(100 * sizeof(int));
	if (data == NULL)
	{
		perror("malloc");
		return 1;
	}

	for (int i = 0; i < 100; i++)
	{
		data[i] = i;
	}

	data[100] = 0;

	printf("data[0]=%d\n", data[0]);
	free(data);
	return 0;
}
\end{verbatim}

\subsection*{直接运行现象}
\begin{verbatim}
data[0]=0
exit:0
\end{verbatim}

\subsection*{valgrind 输出}
\begin{verbatim}
==139266== Memcheck, a memory error detector
==139266== Copyright (C) 2002-2022, and GNU GPL'd, by Julian Seward et al.
==139266== Using Valgrind-3.22.0 and LibVEX; rerun with -h for copyright info
==139266== Command: ./05_oob
==139266== 
==139266== Invalid write of size 4
==139266==    at 0x109216: main (05_oob.c:18)
==139266==  Address 0x4a761d0 is 0 bytes after a block of size 400 alloc'd
==139266==    at 0x4846828: malloc (in /home/amadeus/.local/valgrind/usr/libexec/valgrind/vgpreload_memcheck-amd64-linux.so)
==139266==    by 0x1091BE: main (05_oob.c:6)
==139266== 
data[0]=0
==139266== 
==139266== HEAP SUMMARY:
==139266==     in use at exit: 0 bytes in 0 blocks
==139266==   total heap usage: 2 allocs, 2 frees, 1,424 bytes allocated
==139266== 
==139266== All heap blocks were freed -- no leaks are possible
==139266== 
==139266== For lists of detected and suppressed errors, rerun with: -s
==139266== ERROR SUMMARY: 1 errors from 1 contexts (suppressed: 0 from 0)
\end{verbatim}

\subsection*{分析}
\verb|malloc(100 * sizeof(int))| 仅分配了 100 个 \verb|int| 的空间，有效下标范围为 \verb|0..99|。
语句 \verb|data[100] = 0| 会写入数组末尾之后的 1 个元素位置，属于越界写。

直接运行时程序没有立即崩溃，可能因为越界写不一定立刻触发保护异常，但程序在逻辑上是不正确的。

valgrind 报告 \verb|Invalid write of size 4|，并指出写入地址位于“\verb|0 bytes after a block of size 400 alloc'd|”。
这里 400 字节对应 100 个 \verb|int|，说明越界写刚好发生在分配块的末尾之后，并定位到 \verb|05_oob.c:18|。


\section*{题目6：释放后继续使用（use-after-free）}
\subsection*{源代码}
\begin{verbatim}
#include <stdio.h>
#include <stdlib.h>

int main(void)
{
	int *data = (int *)malloc(100 * sizeof(int));
	if (data == NULL)
	{
		perror("malloc");
		return 1;
	}

	for (int i = 0; i < 100; i++)
	{
		data[i] = i * 2;
	}

	free(data);

	printf("data[10]=%d\n", data[10]);

	return 0;
}
\end{verbatim}

\subsection*{直接运行现象}
\begin{verbatim}
data[10]=20
exit:0
\end{verbatim}

\subsection*{valgrind 输出}
\begin{verbatim}
==139365== Memcheck, a memory error detector
==139365== Copyright (C) 2002-2022, and GNU GPL'd, by Julian Seward et al.
==139365== Using Valgrind-3.22.0 and LibVEX; rerun with -h for copyright info
==139365== Command: ./06_uaf
==139365== 
==139365== Invalid read of size 4
==139365==    at 0x109222: main (06_uaf.c:20)
==139365==  Address 0x4a76068 is 40 bytes inside a block of size 400 free'd
==139365==    at 0x484988F: free (in /home/amadeus/.local/valgrind/usr/libexec/valgrind/vgpreload_memcheck-amd64-linux.so)
==139365==    by 0x109219: main (06_uaf.c:18)
==139365==  Block was alloc'd at
==139365==    at 0x4846828: malloc (in /home/amadeus/.local/valgrind/usr/libexec/valgrind/vgpreload_memcheck-amd64-linux.so)
==139365==    by 0x1091BE: main (06_uaf.c:6)
==139365== 
data[10]=20
==139365== 
==139365== HEAP SUMMARY:
==139365==     in use at exit: 0 bytes in 0 blocks
==139365==   total heap usage: 2 allocs, 2 frees, 1,424 bytes allocated
==139365== 
==139365== All heap blocks were freed -- no leaks are possible
==139365== 
==139365== For lists of detected and suppressed errors, rerun with: -s
==139365== ERROR SUMMARY: 1 errors from 1 contexts (suppressed: 0 from 0)
\end{verbatim}

\subsection*{分析}
在 \verb|free(data)| 之后，\verb|data| 指向的堆块已经被释放；继续读取 \verb|data[10]| 属于释放后使用（Use-After-Free），是未定义行为。

该程序似乎仍然能跑，可能是释放后的堆块短时间内没有被分配器复用/覆盖，所以读到了旧内容。

valgrind 清楚给出 \verb|Invalid read of size 4|，并指明该地址位于“已 free 的 400 字节块内部”，同时给出 malloc/free 的回溯，便于定位问题。


\section*{题目7：向 free 传入异常指针（数组中间位置）}
\subsection*{源代码}
\begin{verbatim}
#include <stdio.h>
#include <stdlib.h>

int main(void)
{
	int *data = (int *)malloc(100 * sizeof(int));
	if (data == NULL)
	{
		perror("malloc");
		return 1;
	}

	for (int i = 0; i < 100; i++)
	{
		data[i] = i;
	}

	free(data + 50);

	printf("unreachable?\n");
	return 0;
}
\end{verbatim}

\subsection*{直接运行现象}
\begin{verbatim}
free(): invalid pointer
exit:134
\end{verbatim}

\subsection*{valgrind 输出}
\begin{verbatim}
==139509== Memcheck, a memory error detector
==139509== Copyright (C) 2002-2022, and GNU GPL'd, by Julian Seward et al.
==139509== Using Valgrind-3.22.0 and LibVEX; rerun with -h for copyright info
==139509== Command: ./07_badfree
==139509== 
==139509== Invalid free() / delete / delete[] / realloc()
==139509==    at 0x484988F: free (in /home/amadeus/.local/valgrind/usr/libexec/valgrind/vgpreload_memcheck-amd64-linux.so)
==139509==    by 0x10921D: main (07_badfree.c:18)
==139509==  Address 0x4a76108 is 200 bytes inside a block of size 400 alloc'd
==139509==    at 0x4846828: malloc (in /home/amadeus/.local/valgrind/usr/libexec/valgrind/vgpreload_memcheck-amd64-linux.so)
==139509==    by 0x1091BE: main (07_badfree.c:6)
==139509== 
unreachable?
==139509== 
==139509== HEAP SUMMARY:
==139509==     in use at exit: 400 bytes in 1 blocks
==139509==   total heap usage: 2 allocs, 2 frees, 1,424 bytes allocated
==139509== 
==139509== 400 bytes in 1 blocks are definitely lost in loss record 1 of 1
==139509==    at 0x4846828: malloc (in /home/amadeus/.local/valgrind/usr/libexec/valgrind/vgpreload_memcheck-amd64-linux.so)
==139509==    by 0x1091BE: main (07_badfree.c:6)
==139509== 
==139509== LEAK SUMMARY:
==139509==    definitely lost: 400 bytes in 1 blocks
==139509==    indirectly lost: 0 bytes in 0 blocks
==139509==      possibly lost: 0 bytes in 0 blocks
==139509==    still reachable: 0 bytes in 0 blocks
==139509==         suppressed: 0 bytes in 0 blocks
==139509== 
==139509== For lists of detected and suppressed errors, rerun with: -s
==139509== ERROR SUMMARY: 2 errors from 2 contexts (suppressed: 0 from 0)
\end{verbatim}

\subsection*{分析}
\verb|free| 的参数必须是动态分配函数（\verb|malloc/calloc/realloc|）返回的“原始指针”（或 \verb|NULL|）。
本题错误地将 \verb|data + 50| 传给 \verb|free|，它指向分配块的中间位置（偏移 50 个 \verb|int|，即 200 字节）。
分配器依赖块头信息管理内存，中间指针无法与正确的块头对应，属于严重的堆管理错误。

直接运行时，glibc 检测到非法指针并报 \verb|free(): invalid pointer|，随后触发 \verb|SIGABRT| 终止（退出码 134）。

valgrind 报告 \verb|Invalid free()|，并指出该地址位于“\verb|200 bytes inside a block of size 400|”，还能定位到 \verb|07_badfree.c:18|。
另外由于没有正确释放原始 \verb|data|，会导致 \verb|definitely lost: 400 bytes| 的泄漏统计。


\section*{题目8：realloc 实现 vector-like 动态数组}
\subsection*{源代码}
\begin{verbatim}
#include <stdio.h>
#include <stdlib.h>

typedef struct
{
	int *arr;
	size_t size;
	size_t cap;
} IntVec;

static void vec_init(IntVec *v)
{
	v->arr = NULL;
	v->size = 0;
	v->cap = 0;
}

static void vec_free(IntVec *v)
{
	free(v->arr);
	v->arr = NULL;
	v->size = 0;
	v->cap = 0;
}

static int vec_push_back(IntVec *v, int x)
{
	if (v->size == v->cap)
	{
		size_t new_cap = (v->cap == 0) ? 1 : (v->cap * 2);
		int *new_arr = (int *)realloc(v->arr, new_cap * sizeof(int));
		if (new_arr == NULL)
		{
			return 0;
		}
		v->arr = new_arr;
		v->cap = new_cap;
	}

	v->arr[v->size] = x;
	v->size++;
	return 1;
}

int main(void)
{
	IntVec v;
	vec_init(&v);

	for (int i = 0; i < 10; i++)
	{
		if (!vec_push_back(&v, i * i))
		{
			fprintf(stderr, "realloc failed\n");
			vec_free(&v);
			return 1;
		}
	}

	printf("size=%zu cap=%zu last=%d\n", v.size, v.cap, v.arr[v.size - 1]);

	vec_free(&v);
	return 0;
}
\end{verbatim}

\subsection*{直接运行现象}
\begin{verbatim}
size=10 cap=16 last=81
exit:0
\end{verbatim}

\subsection*{valgrind 输出}
\begin{verbatim}
==139666== Memcheck, a memory error detector
==139666== Copyright (C) 2002-2022, and GNU GPL'd, by Julian Seward et al.
==139666== Using Valgrind-3.22.0 and LibVEX; rerun with -h for copyright info
==139666== Command: ./08_vector_realloc
==139666== 
size=10 cap=16 last=81
==139666== 
==139666== HEAP SUMMARY:
==139666==     in use at exit: 0 bytes in 0 blocks
==139666==   total heap usage: 6 allocs, 6 frees, 1,148 bytes allocated
==139666== 
==139666== All heap blocks were freed -- no leaks are possible
==139666== 
==139666== For lists of detected and suppressed errors, rerun with: -s
==139666== ERROR SUMMARY: 0 errors from 0 contexts (suppressed: 0 from 0)
\end{verbatim}

\subsection*{性能与对比分析}
该实现使用“容量不足则翻倍”的扩容策略：\verb|cap: 0->1->2->4->8->16...|。
连续插入 $n$ 个元素时，触发 \verb|realloc| 的次数约为 $\log_2 n$；虽然每次扩容可能拷贝已有元素，但总拷贝量为 $1+2+4+\cdots < 2n$，因此 \verb|push_back| 的均摊时间复杂度为 $O(1)$，总体为 $O(n)$。

本程序插入 10 个元素，最终得到 \verb|size=10 cap=16|，符合翻倍扩容的结果。

与链表对比：
\begin{itemize}[leftmargin=2em]
\item vector-like：内存连续，遍历/随机访问 $O(1)$ 且缓存友好；扩容时有一次性搬移成本，但均摊插入仍为 $O(1)$。
\item 链表：若维护尾指针，尾插为 $O(1)$ 且不会整体搬移；但每个节点一次 \verb|malloc/free| 带来较大常数开销，内存碎片多、缓存不友好，随机访问为 $O(n)$。
\end{itemize}

\end{document}